\scanpage{61}
On suppose que dans $\mathcal K$, pour tt diagramme de mono
\diagram{\begin{tikzcd}[ampersand replacement=\&]
\&X\arrow[dl,hook,"{i}"']\arrow[dr,hook,"{j}"]\&\\
Y\&\&Z
\end{tikzcd}}
il existe une somme amalgamée $S=Y\amalg_X Z$, et que de plus $Y,Z\subseteq S$ sont des mono et $Y\times_S Z\overset{\sim}{\longleftarrow}X$.
\unsure{Il revient} au $=$ de dire \deleted{(par le ``de plus'')} \inserted{et encore}
(compte tenu des hyp \deleted{de} \inserted{d'existence} \unsure{cf. aussi ci-dessus} faites plus haut sur $\mathcal K$) qu'il existe un diagramme commut \inserted{\uline{cartésien}} de mono
\diagram{(D)\quad\begin{tikzcd}[ampersand replacement=\&]
\&X\arrow[dl,hook,"{i}"']\arrow[dr,hook,"{j}"]\&\\
Y\arrow[dr,hook]\&\&Z\arrow[dl,hook]\\
\&S'\&
\end{tikzcd}}
(et alors on peut prendre $S=Y\cup_{S'}Z$ (sup de sous-objets).
Il résulte de ceci que tout mono de $\mathcal K$ est effectif (peut être pas universel---).
Les sommes amalgamées sont alors universelles (stables par chgt de base quelconque).

Supposons maintenant que $X,Y,Z$ et $i,j$ soient modérés. On cherche sur $S$ une structure modérée telle que $Y-S$ et $Z-S$ soient modérés (auquel cas ce sera une somme amalgamée \uline{dans} $\mathcal M$).
\unsure{Cela|Celle} existe ssi on peut trouver un diagramme \uline{cartésien} (D) avec $S'$ et $Y,Z-S'$ \uline{modérés}.

Ça marchera si on suppose qu'on a de $=$
\inserted{$0_1:e_0-I_0$ dans $\mathcal M$ (cf $4^\circ$)}
\margin{d'où $0_n:e\longrightarrow I_0^n$}
\[
\left\{\begin{array}{ll}
5^\circ)& I_0\text{ injectif dans }\mathcal M\text{ i.e. }\forall\text{ mono }Y\lhook\joinrel\xrightarrow{i}X\text{ dans}\\
&\mathcal M,\text{ et tt }f:Y-I_0,\ f\text{ se prolonge à }X\quad(\text{NB. OPS }X=I_0^n)\\
6^\circ)&\forall\text{ mono }Y\lhook\joinrel\longrightarrow X,\ \deleted{\text{dans}}\ \exists X\xrightarrow{f}I_0^n\text{ tel que}\\
&Y=f^{-1}(0_n),\text{ i.e. diag. cartésien}
\end{array}\right.
\]
\diagram{\begin{tikzcd}[ampersand replacement=\&]
X\arrow[r,dash]\&I_0^n\\
Y\arrow[u,hook]\arrow[r]\&e\arrow[u,"{0_n}"']
\end{tikzcd}}
(NB OPS $X=I_0^n$)

\margin{62}
